\chapter{Closing the \texorpdfstring{$\kappa$}{kappa}-curve: the reduction is justified}
\label{chap:reduction}

% archon:covers Gluck/Euclidean/Reduction.lean

This chapter assembles the winding-number argument (\cref{chap:winding}) and the
preliminary reparametrization (\cref{chap:stepreduction}) into the load-bearing
analytic statement \cref{lem:reduction_justified}: every four-vertex curvature
function $\kappa$ admits a circle reparametrization $h$ for which the
reconstruction weight $1/(\kappa\circ h)$ has vanishing error vector, so the
reconstruction curve closes up.

\medskip

\noindent\emph{Why a new chapter / file.} The conclusion consumes
\cref{thm:existence_of_zero} from \cref{chap:winding}, which imports
\cref{chap:bicircle} and hence \cref{chap:stepreduction}; so the statement cannot
live in \texttt{StepReduction.lean} (that would be an import cycle). It is placed
in a new file \texttt{Reduction.lean} that \texttt{import}s
\texttt{Gluck.Winding} (transitively pulling in everything else).

\medskip

\noindent\emph{The genuine remaining construction.} \cref{thm:existence_of_zero}
produces a zero of the \emph{bicircle} (step-function) error map
$\mathrm{errorMap} = E_{\mathrm{bi}}$ --- a closing configuration of the
two-valued step family --- but \emph{not} a curve realising the original
continuous $\kappa$. To transfer the closure to $\kappa$ we re-run the planar
degree principle on a \emph{$\kappa$-error map} $E_\kappa$ over the same disk,
obtained by feeding $\kappa\circ h_1$ through a family of breakpoint-aligning
reparametrizations $g_z$, and showing $E_\kappa$ is uniformly close to
$E_{\mathrm{bi}}$ on $\partial\mathbb{D}$ (so it inherits the nonzero boundary
winding). The robustness is exactly the $L^1$/measure continuity of the error
vector (\cref{lem:error_vector_lipschitz}), \emph{not} the false ``$C^1$-close
curves'' claim the survey uses (see \cref{chap:stepreduction}'s departure note).

\section{The breakpoint-aligning reparametrization family \texorpdfstring{$g_z$}{g\_z}}

Recall the centred bicircle configuration $\mathrm{configSpace}(\delta)(z) =
(\tfrac{\pi}{4}+\delta\,z.\mathrm{re},\ \tfrac{3\pi}{4}+\delta\,z.\mathrm{im},\
\tfrac{5\pi}{4},\ \tfrac{7\pi}{4})$ (\cref{def:configuration_space}): only the two
leading breakpoints move with $z\in\overline{\mathbb{D}}$, the trailing two are
pinned. The preliminary reparametrization $h_1$ of
\cref{lem:exists_preliminary_reparam} makes $\kappa\circ h_1$ $\varepsilon$-close
in measure to the canonical four-arc step function $\kappa_0 = \mathrm{stepCurvature}\,
b\,a\,0\,\tfrac{\pi}{2}\,\pi\,\tfrac{3\pi}{2}$ (breakpoints $0,\tfrac{\pi}{2},
\pi,\tfrac{3\pi}{2}$, values $a,b,a,b$).

The four arcs
$[\theta_1,\theta_2),[\theta_2,\theta_3),[\theta_3,\theta_4),[\theta_4,\theta_1+2\pi)$
have lengths
\[
  L_1=\tfrac{\pi}{2}+\delta(z.\mathrm{im}-z.\mathrm{re}),\quad
  L_2=\tfrac{\pi}{2}-\delta z.\mathrm{im},\quad
  L_3=\tfrac{\pi}{2},\quad
  L_4=\tfrac{\pi}{2}+\delta z.\mathrm{re};
\]
for $\delta\le\pi/8$ and $z\in\overline{\mathbb D}$ each $L_k$ lies in
$[\tfrac{\pi}{4},\tfrac{3\pi}{4}]$ and depends continuously on $z$.

\medskip

\noindent\emph{Construction by a calibrated density (replaces the discarded
piecewise-linear/tent recipe).} We build $g_z$ as the running integral of a
\emph{calibrated speed density} $w_z$ rather than as a piecewise-linear map. The
running-integral form is what makes continuity, strict monotonicity, quasi-period\-icity
and joint continuity in $(z,\theta)$ clean to prove (FTC and parametric-integral
continuity), and it avoids the $\mathrm{toIcoMod}$-periodization of a piecewise-linear
map, which has an intractable seam/joint-continuity problem. (Historical note: a
finite sum of periodic tents $\arccos(\cos(\theta-c))$ does \emph{not} produce an
exact four-corner piecewise-linear \emph{node-mapping} --- each such tent is a
symmetric triangle wave with two corners per period, so four tents give eight
corners, the spurious four corrupting the node mapping. This is why $g_z$ is built
from a density rather than directly as a tent-sum. Note this is no contradiction
with \cref{chap:stepreduction}, where $\arccos(\cos)$ tents are used \emph{correctly}
as continuous building blocks for a continuous \emph{density} --- a use for which
their two-corner symmetric shape is harmless; the obstruction is specific to forcing
exact corner positions on the map itself.)

\begin{definition}
[Calibrated speed density $w_z$]
  \label{def:align_density}
  \lean{Gluck.clampTent}
  % Lean (private, not pinned): alignDensity, arccos_cos_abs, alignN1, alignN2, alignN3, alignN4, alignL1, alignL2, alignL3, alignL4, alignC1, alignC2, alignC3, alignC4, alignHt, alignHt_eq
  \uses{def:configuration_space}
  Project-bespoke. Fix the ramp half-width $\eta=\pi/16$ (so $2\eta=\pi/8$ is
  strictly less than the minimal arc length $\pi/4$, hence the four node windows
  $[\theta_k-\eta,\theta_k+\eta]$ are pairwise disjoint and each arc retains a
  non-empty central plateau of length $L_k-2\eta\ge\pi/8$). Let $w_z:\mathbb R\to\mathbb R$
  be the $2\pi$-periodic, continuous, piecewise-affine
  (\emph{trapezoidal}) density determined by:
  \begin{itemize}
    \item a common node value $w_z(\theta_k)=V$ at all four nodes, where
      $V:=2/3$ (a fixed positive constant);
    \item a constant plateau value $m_k(z)$ on the central part
      $[\theta_k+\eta,\theta_{k+1}-\eta]$ of arc $k$;
    \item linear ramps on each node window $[\theta_k-\eta,\theta_k+\eta]$
      joining the plateau value $m_{k-1}(z)$ (left) to the plateau value
      $m_k(z)$ (right) through the common node value $V$ at the centre
      $\theta_k$ --- concretely $w_z$ is affine from $m_{k-1}(z)$ at
      $\theta_k-\eta$ to $V$ at $\theta_k$ and affine from $V$ at $\theta_k$ to
      $m_k(z)$ at $\theta_k+\eta$.
  \end{itemize}
  The plateau height $m_k(z)$ is solved so that the arc integral is exactly
  $\pi/2$: with $\int_{\theta_k}^{\theta_{k+1}}w_z =
  m_k(z)\,(L_k-2\eta) + \tfrac{\eta}{2}(V+m_k(z)) + \tfrac{\eta}{2}(m_k(z)+V)
  = m_k(z)(L_k-\eta) + \eta V$, setting this to $\pi/2$ gives
  \[
    m_k(z) \;=\; \frac{\pi/2-\eta V}{L_k-\eta}.
  \]
  Since $L_k\in[\tfrac\pi4,\tfrac{3\pi}4]$, $\eta=\pi/16$ and $V=2/3$, the numerator
  $\pi/2-\eta V = \pi/2-\pi/24>0$ and the denominator $L_k-\eta\in[\tfrac{3\pi}{16},\tfrac{11\pi}{16}]$,
  so each $m_k(z)$ is a continuous, strictly positive function of $z$. Since $w_z$
  is a convex combination of the node value $V$ and the plateau heights $m_k(z)$
  (it is affine between them on the ramps), $w_z(\theta)\in[c_-,c_+]\subset(0,\infty)$
  with $c_-=\min(V,\min_k m_k)$ and $c_+=\max(V,\max_k m_k)$. A direct computation
  over $L_k\in[\tfrac\pi4,\tfrac{3\pi}4]$ gives $\min_k m_k=2/3$ (at $L_k=3\pi/4$)
  and $\max_k m_k=22/9$ (at $L_k=\pi/4$), so $c_-=2/3$ and $c_+=22/9$; we set the
  uniform slope bound $K:=\max(c_+,1/c_-)=22/9$ (so $1/K\le w_z\le K$). This is the
  same kind of explicit continuous trapezoidal-density construction already carried
  out (for a different calibration) in the preliminary reparametrization of
  \cref{chap:stepreduction}; the present density is self-contained and depends on
  none of that chapter's lemmas.

  \emph{Lean realization.} Each ramp/plateau pulse is the periodic trapezoid
  $\mathrm{clampTent}(\eta,L,\tau)(\theta)=\min(1,\max(0,(L/2-\arccos\cos(\theta-\tau))/\eta))$
  (height $1$, centre $\tau$, width $L$, ramp $\eta$), built on the
  $\arccos\cos$ idiom of \cref{chap:stepreduction} for free $2\pi$-periodicity and
  continuity; $w_z=\mathrm{alignDensity}$ is the constant $V$ plus a $z$-continuous
  combination of four such pulses with heights $m_k(z)=\mathrm{alignHt}(L_k)$ at the
  four nodes $\mathrm{alignN}_k(z)$ / arc lengths $\mathrm{alignL}_k(z)$ / centres
  $\mathrm{alignC}_k(z)$.  These named pieces (and their continuity helpers) carry the
  dependency edges; their properties are collected in \cref{lem:align_density_props}
  and \cref{lem:align_density_arc_integral}.
\end{definition}

\begin{lemma}
[Density is well-behaved]
  \label{lem:align_density_props}
  \lean{Gluck.continuous_clampTent, Gluck.continuous_clampTent_theta, Gluck.clampTent_nonneg, Gluck.clampTent_le_one, Gluck.clampTent_periodic}
  % Lean (private, not pinned): continuous_uncurry_alignDensity, continuous_alignDensity_theta, alignDensity_periodic, alignDensity_ge, alignHt_bounds, alignL_bounds, continuous_alignHt, continuous_alignN1, continuous_alignN2, continuous_alignN3, continuous_alignN4, continuous_alignL1, continuous_alignL2, continuous_alignL3, continuous_alignL4, continuous_alignC1, continuous_alignC2, continuous_alignC3, continuous_alignC4, alignDensity_le, alignHt_le, alignHt_nonneg
  \uses{def:align_density}
  For $0<\delta\le\pi/8$ and $z\in\overline{\mathbb D}$: $w_z$ is continuous,
  $2\pi$-periodic, the joint map $(z,\theta)\mapsto w_z(\theta)$ is continuous on
  $\overline{\mathbb D}\times\mathbb R$, and $c_-\le w_z(\theta)\le c_+$ for
  absolute constants $0<c_-\le c_+$.
\end{lemma}
\begin{proof}
  \uses{def:align_density}
  Continuity and periodicity are immediate from the trapezoidal piecewise-affine
  shape (the pieces agree at the breakpoints by the common node value $V$ and the
  shared plateau values). Joint continuity holds because the node positions
  $\theta_k(z)$ and the plateau heights $m_k(z)$ are continuous in $z$ (the arc
  lengths $L_k$ stay in $[\tfrac\pi4,\tfrac{3\pi}4]$, bounded away from the
  $L_k=\eta$ singularity), and an affine interpolation with continuously varying
  knots and node values is jointly continuous. The bound is the height estimate of
  \cref{def:align_density}.
\end{proof}

\begin{lemma}

\leanok
[Disk smallness of the breakpoint shifts]
  \label{lem:align_delta_bounds}
  % Lean (private, not pinned): align_delta_bounds
  \uses{def:align_density}
  For $0<\delta\le\pi/8$ and $\|z\|\le1$, the two varying breakpoint shifts satisfy
  $\delta\,z.\mathrm{re},\ \delta\,z.\mathrm{im}\in[-\tfrac\pi8,\tfrac\pi8]$.
\end{lemma}
\begin{proof}
  \uses{def:align_density}
  \leanok
  $|z.\mathrm{re}|,\,|z.\mathrm{im}|\le\|z\|\le1$, so
  $|\delta\,z.\mathrm{re}|,\,|\delta\,z.\mathrm{im}|\le\delta\le\pi/8$.
\end{proof}

\begin{lemma}
[Exact arc integrals]
  \label{lem:align_density_arc_integral}
  \lean{Gluck.clampTent_integral_support, Gluck.clampTent_eq_zero, Gluck.half_le_arccos_cos}
  % Lean (private, not pinned): alignDensity_arc1, alignDensity_arc2, alignDensity_arc3, alignDensity_arc4, alignDensity_period_integral, alignDensity_integral_split, clampTent_centered_integral, clampTent_integral_eq_zero
  \uses{def:align_density, lem:align_delta_bounds}
  For each arc $k\in\{1,2,3,4\}$, $\int_{\theta_k}^{\theta_{k+1}}w_z=\tfrac\pi2$,
  and hence the full-period integral $\int_{\theta_1}^{\theta_1+2\pi}w_z=2\pi$.
\end{lemma}
\begin{proof}
  \uses{def:align_density}
  The arc integral splits into the right half of the node-$\theta_k$ ramp
  ($\int=\tfrac\eta2(V+m_k)$), the central plateau ($\int=(L_k-2\eta)m_k$), and the
  left half of the node-$\theta_{k+1}$ ramp ($\int=\tfrac\eta2(m_k+V)$); the sum is
  $m_k(L_k-\eta)+\eta V$, which equals $\pi/2$ by the choice of $m_k$ in
  \cref{def:align_density}. Summing the four arcs gives $4\cdot\tfrac\pi2=2\pi$.
\end{proof}

\section{The generic running-integral reparametrization}

The structural facts about $g_z$ --- the fundamental theorem of calculus,
continuity, strict monotonicity, quasi-periodicity, and the one-period change of
variables --- do not depend on the specific calibrated density $w_z$; they hold for
the running integral of any continuous (and, where needed, strictly positive,
$2\pi$-periodic) weight. We isolate them once on a generic family
$\mathrm{integralReparam}$ and let the $g_z$ lemmas delegate to them.

\begin{definition}

\leanok
[Running-integral reparametrization]
  \label{def:integral_reparam}
  \lean{Gluck.integralReparam}
  Project-bespoke. For a weight $w:\mathbb R\to\mathbb R$ and an anchor constant
  $c\in\mathbb R$, define
  \[
    \mathrm{integralReparam}\,w\,c \;:\; s \;\longmapsto\; c + \int_0^{s} w(t)\,dt .
  \]
\end{definition}

\begin{lemma}

\leanok
[FTC derivative of the running integral]
  \label{lem:has_deriv_integral_reparam}
  \lean{Gluck.hasDerivAt_integralReparam}
  \uses{def:integral_reparam}
  If $w$ is continuous, then $\mathrm{integralReparam}\,w\,c$ is differentiable with
  $(\mathrm{integralReparam}\,w\,c)'(s) = w(s)$ for every $s$.
\end{lemma}
\begin{proof}
  \uses{def:integral_reparam}
  \leanok
  By the fundamental theorem of calculus the derivative of the primitive
  $s\mapsto\int_0^{s}w$ of the continuous integrand $w$ is $w(s)$; adding the
  constant $c$ does not change the derivative.
\end{proof}

\begin{lemma}

\leanok
[Continuity of the running integral]
  \label{lem:continuous_integral_reparam}
  \lean{Gluck.continuous_integralReparam}
  \uses{def:integral_reparam}
  If $w$ is continuous, then $\mathrm{integralReparam}\,w\,c$ is continuous.
\end{lemma}
\begin{proof}
  \uses{def:integral_reparam}
  \leanok
  The primitive of a continuous (hence locally interval-integrable) function is
  continuous, and adding the constant $c$ preserves continuity.
\end{proof}

\begin{lemma}

\leanok
[Strict monotonicity from a positive density]
  \label{lem:strict_mono_integral_reparam}
  \lean{Gluck.strictMono_integralReparam}
  \uses{def:integral_reparam}
  If $w$ is continuous and strictly positive everywhere ($0<w(s)$ for all $s$), then
  $\mathrm{integralReparam}\,w\,c$ is strictly increasing.
\end{lemma}
\begin{proof}
  \uses{def:integral_reparam}
  \leanok
  For $x<y$,
  $\mathrm{integralReparam}\,w\,c\,(y)-\mathrm{integralReparam}\,w\,c\,(x)
  =\int_x^y w>0$, since the integrand is strictly positive on the nondegenerate
  interval $[x,y]$.
\end{proof}

\begin{lemma}

\leanok
[Quasi-periodicity of the running integral]
  \label{lem:integral_reparam_add_two_pi}
  \lean{Gluck.integralReparam_add_two_pi}
  \uses{def:integral_reparam}
  If $w$ is continuous, $2\pi$-periodic, with full-period integral
  $\int_0^{2\pi} w = 2\pi$, then
  $\mathrm{integralReparam}\,w\,c\,(s+2\pi)=\mathrm{integralReparam}\,w\,c\,(s)+2\pi$.
\end{lemma}
\begin{proof}
  \uses{def:integral_reparam}
  \leanok
  Split $\int_0^{s+2\pi}=\int_0^{s}+\int_{s}^{s+2\pi}$; by $2\pi$-periodicity of $w$
  the second integral equals the full-period integral $\int_0^{2\pi} w = 2\pi$.
\end{proof}

\begin{lemma}

\leanok
[Change of variables for the running integral]
  \label{lem:integral_reparam_change_of_var}
  \lean{Gluck.integralReparam_changeOfVar}
  \uses{def:integral_reparam, lem:has_deriv_integral_reparam,
        lem:continuous_integral_reparam, lem:strict_mono_integral_reparam}
  For a continuous, strictly-positive weight $w$ and any integrand $G$, writing
  $g=\mathrm{integralReparam}\,w\,c$,
  \[
    \int_{g(0)}^{g(2\pi)} G(x)\,dx \;=\; \int_0^{2\pi} w(x)\,G\bigl(g(x)\bigr)\,dx .
  \]
\end{lemma}
\begin{proof}
  \uses{lem:has_deriv_integral_reparam, lem:continuous_integral_reparam,
        lem:strict_mono_integral_reparam}
        \leanok
  $g$ is strictly monotone (\cref{lem:strict_mono_integral_reparam}), hence
  injective, and $C^1$ with derivative $w$ (\cref{lem:has_deriv_integral_reparam});
  being continuous (\cref{lem:continuous_integral_reparam}) and monotone it maps
  $[0,2\pi]$ onto $[g(0),g(2\pi)]$. The one-dimensional substitution rule
  $\int_{g(I)}G=\int_I |g'|\,(G\circ g)$ applies, and $|g'|=w>0$ drops the absolute
  value.
\end{proof}

\begin{definition}

\leanok
[Breakpoint-aligning reparametrization family]
  \label{def:align_reparam}
  % Lean (private, not pinned): alignReparam
  \uses{def:align_density, def:configuration_space, def:four_arc_step_function}
  Project-bespoke. For $0<\delta\le\pi/8$ and $z\in\overline{\mathbb{D}}$ define
  $g_z := \mathrm{alignReparam}(\delta)(z) : \mathbb{R}\to\mathbb{R}$ as the running
  integral of the calibrated density $w_z$ (\cref{def:align_density}) anchored so
  that $g_z(\theta_1)=\pi/2$:
  \[
    g_z(\theta) \;:=\; \tfrac{\pi}{2} + \int_{\theta_1}^{\theta} w_z(t)\,dt .
  \]
  $g_z$ is the orientation-preserving circle homeomorphism that sends the four
  breakpoints of the $z$-configuration
  \[
    \theta_1=\tfrac{\pi}{4}+\delta z.\mathrm{re},\quad
    \theta_2=\tfrac{3\pi}{4}+\delta z.\mathrm{im},\quad
    \theta_3=\tfrac{5\pi}{4},\quad
    \theta_4=\tfrac{7\pi}{4}
  \]
  to the canonical breakpoints $\tfrac{\pi}{2},\pi,\tfrac{3\pi}{2},2\pi$
  respectively (\cref{lem:align_reparam_node_values}), is strictly increasing with
  slope $g_z'=w_z\in[c_-,c_+]\subset(0,\infty)$, and is quasi-periodic
  ($g_z(\theta+2\pi)=g_z(\theta)+2\pi$) because the full-period integral of $w_z$ is
  $2\pi$ (\cref{lem:align_density_arc_integral}). We write $K=22/9$ for the uniform
  slope bound ($1/K\le w_z\le K$; see \cref{def:align_density}).
\end{definition}

\begin{lemma}

\leanok
[$g_z$ is an instance of the generic family]
  \label{lem:align_reparam_eq}
  % Lean (private, not pinned): alignReparam_eq
  \uses{def:align_reparam, def:integral_reparam, def:align_density}
  For every $\delta,z$,
  \[
    \mathrm{alignReparam}(\delta)(z)
      = \mathrm{integralReparam}\,(w_z)\,\Bigl(\tfrac\pi2-\int_0^{\theta_1} w_z\Bigr),
  \]
  where $\theta_1=\mathrm{alignN1}(\delta)(z)$ is the first configuration breakpoint.
\end{lemma}
\begin{proof}
  \uses{def:align_reparam, def:integral_reparam, def:align_density}
  \leanok
  By definition $g_z(\theta)=\tfrac\pi2+\int_{\theta_1}^{\theta} w_z$, and
  $\int_{\theta_1}^{\theta} w_z = \int_0^{\theta} w_z - \int_0^{\theta_1} w_z$, so the
  anchor $\theta_1$ folds into the additive constant
  $\tfrac\pi2-\int_0^{\theta_1} w_z$, exhibiting $g_z$ as
  $\mathrm{integralReparam}\,(w_z)$ at that anchor.
\end{proof}

These structural properties are recorded as separate Lean lemmas (used by
\cref{lem:kappa_error_map_continuous} and \cref{lem:kappa_error_map_close}):

\begin{lemma}

\leanok
[Node values of $g_z$]
  \label{lem:align_reparam_node_values}
  % Lean (private, not pinned): alignReparam_node_values
  \uses{def:align_reparam, lem:align_density_arc_integral}
  $g_z(\theta_1)=\tfrac\pi2$, $g_z(\theta_2)=\pi$, $g_z(\theta_3)=\tfrac{3\pi}2$,
  $g_z(\theta_4)=2\pi$; i.e.\ $g_z(\theta_k)=k\cdot\tfrac\pi2$.
\end{lemma}
\begin{proof}
  \uses{def:align_reparam, lem:align_density_arc_integral}
  \leanok
  $g_z(\theta_1)=\pi/2$ by definition. By
  \cref{lem:align_density_arc_integral}, $g_z(\theta_{k+1})-g_z(\theta_k)=
  \int_{\theta_k}^{\theta_{k+1}}w_z=\pi/2$, so $g_z(\theta_k)=\pi/2+(k-1)\pi/2=k\pi/2$
  by induction.
\end{proof}

\begin{lemma}

\leanok
[Quasi-periodicity of $g_z$]
  \label{lem:align_reparam_add_two_pi}
  % Lean (private, not pinned): alignReparam_add_two_pi
  \uses{def:align_reparam, lem:align_density_arc_integral,
        lem:integral_reparam_add_two_pi, lem:align_reparam_eq}
  $g_z(\theta+2\pi)=g_z(\theta)+2\pi$ for all $\theta$.
\end{lemma}
\begin{proof}
  \uses{def:align_reparam, lem:align_density_arc_integral}
  \leanok
  $g_z(\theta+2\pi)-g_z(\theta)=\int_{\theta}^{\theta+2\pi}w_z=
  \int_{\theta_1}^{\theta_1+2\pi}w_z=2\pi$ by the $2\pi$-periodicity of $w_z$ (the
  integral of a periodic function over any window of length one period is the
  full-period integral) and \cref{lem:align_density_arc_integral}.
\end{proof}

\begin{lemma}

\leanok
[Continuity of $g_z$ in $\theta$]
  \label{lem:continuous_align_reparam}
  % Lean (private, not pinned): continuous_alignReparam
  \uses{def:align_reparam, lem:align_density_props,
        lem:continuous_integral_reparam, lem:align_reparam_eq}
  $g_z$ is continuous for each fixed $\delta,z$.
\end{lemma}
\begin{proof}
  \uses{def:align_reparam, lem:align_density_props}
  \leanok
  The running integral $\theta\mapsto\int_{\theta_1}^{\theta}w_z$ of a continuous
  integrand is continuous (indeed $C^1$ by FTC, with derivative $w_z$).
\end{proof}

\begin{lemma}

\leanok
[Joint continuity of $(z,\theta)\mapsto g_z(\theta)$]
  \label{lem:continuous_uncurry_align_reparam}
  % Lean (private, not pinned): continuous_uncurry_alignReparam
  \uses{def:align_reparam, lem:align_density_props}
  The map $(z,\theta)\mapsto g_z(\theta)$ is jointly continuous on
  $\overline{\mathbb D}\times\mathbb R$. This is the load-bearing input to
  \cref{lem:kappa_error_map_continuous}.
\end{lemma}
\begin{proof}
  \uses{def:align_reparam, lem:align_density_props}
  \leanok
  The integrand $(z,t)\mapsto w_z(t)$ is jointly continuous
  (\cref{lem:align_density_props}) and the lower limit $\theta_1(z)$ is continuous
  in $z$, so continuity of the parametrised integral with a jointly continuous
  integrand over a varying interval (the same parametric-integral continuity used
  for \cref{lem:kappa_error_map_continuous}) gives joint continuity of
  $(z,\theta)\mapsto g_z(\theta)$.
\end{proof}

\begin{lemma}

\leanok
[Strict monotonicity of $g_z$ on the disk]
  \label{lem:strict_mono_align_reparam}
  % Lean (private, not pinned): strictMono_alignReparam
  \uses{def:align_reparam, lem:align_density_props,
        lem:strict_mono_integral_reparam, lem:align_reparam_eq}
  For $0<\delta\le\pi/8$ and $\|z\|\le1$, $g_z$ is strictly increasing.
\end{lemma}
\begin{proof}
  \uses{def:align_reparam, lem:align_density_props}
  \leanok
  $g_z(y)-g_z(x)=\int_x^y w_z\ge c_-\,(y-x)>0$ for $x<y$, since the integrand
  $w_z\ge c_->0$ (\cref{lem:align_density_props}).
\end{proof}

\begin{lemma}

\leanok
[$g_z$ aligns the step functions]
  \label{lem:align_reparam_matches}
  % Lean (private, not pinned): kappaZero_comp_alignReparam
  \uses{def:align_reparam, lem:align_reparam_node_values,
        lem:strict_mono_align_reparam, def:four_arc_step_function,
        def:configuration_space}
  Let $\kappa_0=\mathrm{stepCurvature}\,b\,a\,0\,\tfrac\pi2\,\pi\,\tfrac{3\pi}2$
  be the canonical four-arc step function of
  \cref{lem:exists_preliminary_reparam}. Then for almost every $\theta$,
  \[
    \kappa_0\bigl(g_z(\theta)\bigr)
      \;=\; \mathrm{stepCurvature}\,a\,b\,(\mathrm{configSpace}(\delta)(z))\,(\theta),
  \]
  the step function of the $z$-configuration. Consequently
  $\mathrm{errorMap}(a,b,\delta)(z)=\mathrm{errorVector}(\mathrm{radius}(\kappa_0\circ g_z))$.
\end{lemma}
\begin{proof}
  \uses{def:align_reparam, lem:align_reparam_node_values, lem:strict_mono_align_reparam}
  \leanok
  By \cref{lem:align_reparam_node_values} and strict monotonicity
  (\cref{lem:strict_mono_align_reparam}),
  $g_z$ carries arc~1 $[\theta_1,\theta_2)\to[\tfrac\pi2,\pi)$, arc~2
  $\to[\pi,\tfrac{3\pi}2)$, arc~3 $\to[\tfrac{3\pi}2,2\pi)$, arc~4
  $\to[2\pi,\tfrac\pi2+2\pi)\equiv[0,\tfrac\pi2)$ (mod $2\pi$). On these target
  arcs $\kappa_0$ takes the values $b,a,b,a$ respectively, which match the values
  $b,a,b,a$ of $\mathrm{stepCurvature}\,a\,b\,(\mathrm{configSpace})$ on arcs
  $1,2,3,4$. The two sides therefore agree off the finite breakpoint set (a null
  set). Since $\mathrm{errorMap}(a,b,\delta)(z)=
  \mathrm{bicircleErrorVector}\,a\,b\,(\mathrm{configSpace})=
  \mathrm{errorVector}(\mathrm{radius}(\mathrm{stepCurvature}\,a\,b\,(\mathrm{configSpace})))$
  and $\mathrm{errorVector}$ depends only on the a.e.\ class of the weight, the
  stated identity of error vectors follows.
\end{proof}

\section{The \texorpdfstring{$\kappa$}{kappa}-error map over the disk}

\begin{definition}


\leanok
[The $\kappa$-error map]
  \label{def:kappa_error_map}
  % Lean (private, not pinned): kappaErrorMap
  \uses{def:align_reparam, def:error_vector, def:radius}
  Project-bespoke. For a curvature function $\kappa$, the reparametrization $h_1$,
  and $0<\delta\le\pi/8$,
  \[
    E_\kappa(z) \;:=\; \mathrm{errorVector}\bigl(\mathrm{radius}(\kappa\circ h_1\circ g_z)\bigr)
    \;=\; \int_0^{2\pi} e^{i\theta}\,\frac{d\theta}{\kappa\bigl(h_1(g_z(\theta))\bigr)},
    \qquad z\in\overline{\mathbb{D}} .
  \]
  This is the error vector of the curve reconstructed from the curvature function
  $\kappa\circ h_1\circ g_z$; $E_\kappa(z)=0$ means precisely that this curve
  closes up.
\end{definition}

\begin{lemma}


\leanok
[$E_\kappa$ is continuous on the disk]
  \label{lem:kappa_error_map_continuous}
  % Lean (private, not pinned): continuous_kappaErrorMap
  \uses{def:kappa_error_map, def:align_reparam}
  $E_\kappa : \overline{\mathbb{D}}\to\mathbb{C}$ is continuous.
\end{lemma}
\begin{proof}
  \uses{def:kappa_error_map, def:align_reparam}
  \leanok
  The integrand $(z,\theta)\mapsto e^{i\theta}/\kappa(h_1(g_z(\theta)))$ is jointly
  continuous on $\overline{\mathbb{D}}\times[0,2\pi]$: $\kappa,h_1$ are continuous,
  $(z,\theta)\mapsto g_z(\theta)$ is jointly continuous (\cref{def:align_reparam}),
  and $\kappa>0$ keeps the denominator bounded away from $0$ on the compact domain.
  Continuity of a parametrised interval integral with jointly continuous integrand
  over a fixed compact interval gives continuity of $z\mapsto E_\kappa(z)$.
\end{proof}

\section{Change-of-variables keystones for the \texorpdfstring{$L^1$}{L1} estimate}

The closeness estimate \cref{lem:kappa_error_map_close} and its existence form
\cref{lem:exists_reparam_close} both pull an integral over $\theta$ back to the
$\phi=g_z(\theta)$ line through the slope-bounded reparametrization $g_z$. Two
analytic facts about $g_z$ carry this; both are proved.

\begin{lemma}

\leanok
[FTC derivative of $g_z$]
  \label{lem:has_deriv_align_reparam}
  % Lean (private, not pinned): hasDerivAt_alignReparam
  \uses{def:align_reparam, lem:align_density_props,
        lem:has_deriv_integral_reparam, lem:align_reparam_eq}
  $g_z$ is differentiable with $g_z'(\theta)=w_z(\theta)$ for every $\theta$.
\end{lemma}
\begin{proof}
  \uses{def:align_reparam, lem:align_density_props}
  \leanok
  The running integral $\theta\mapsto\int_{\theta_1}^{\theta}w_z$ of the continuous
  integrand $w_z$ has derivative $w_z(\theta)$ at every point by the fundamental
  theorem of calculus (the right-endpoint derivative of an interval integral of a
  continuous function).
\end{proof}

\begin{lemma}

\leanok
[Change of variables through $g_z$]
  \label{lem:align_reparam_change_of_var}
  \lean{Gluck.alignReparam_changeOfVar}
  \uses{lem:has_deriv_align_reparam, lem:strict_mono_align_reparam,
        lem:integral_reparam_change_of_var, lem:align_reparam_eq}
  For any measurable $G$,
  $\int_{[g_z(0),\,g_z(2\pi)]} G
     = \int_{[0,2\pi]} w_z(\theta)\,G(g_z(\theta))\,d\theta$.
\end{lemma}
\begin{proof}
  \uses{lem:has_deriv_align_reparam, lem:strict_mono_align_reparam}
  \leanok
  $g_z$ is strictly monotone (\cref{lem:strict_mono_align_reparam}), hence injective,
  and $C^1$ with derivative $w_z\ge 0$ (\cref{lem:has_deriv_align_reparam}). The
  image of $[0,2\pi]$ is $[g_z(0),g_z(2\pi)]$ (continuous monotone image of an
  interval), so the one-dimensional change-of-variables formula
  $\int_{g(S)}G=\int_S |g'|\,(G\circ g)$ applies with $|g'|=w_z$.
\end{proof}

\section{Uniform closeness to the bicircle error map}

\begin{lemma}


\leanok
[$E_\kappa$ is uniformly close to $E_{\mathrm{bi}}$]
  \label{lem:kappa_error_map_close}
  % Lean (private, not pinned): kappaErrorMap_sub_errorMap_le
  \uses{def:kappa_error_map, def:error_map, lem:error_vector_lipschitz,
        lem:exists_preliminary_reparam, def:align_reparam,
        lem:align_reparam_matches, lem:strict_mono_align_reparam,
        lem:has_deriv_align_reparam, lem:align_reparam_change_of_var}
  Let $M = 1/\min\kappa$ and let $h_1$ be chosen for a tolerance $\varepsilon>0$
  as in \cref{lem:exists_preliminary_reparam}. There is an absolute constant
  $K$ (the slope bound of $g_z$) such that, for every $z\in\overline{\mathbb{D}}$,
  \[
    \bigl\lVert E_\kappa(z) - \mathrm{errorMap}(a,b,\delta)(z)\bigr\rVert
      \;\le\; K\,M^2\,\varepsilon .
  \]
\end{lemma}
\begin{proof}
  \uses{lem:error_vector_lipschitz, lem:exists_preliminary_reparam, def:align_reparam,
        lem:align_reparam_matches, lem:strict_mono_align_reparam}
        \leanok
  Fix $z$. By \cref{lem:align_reparam_matches},
  $\mathrm{errorMap}(a,b,\delta)(z) = \mathrm{errorVector}(\mathrm{radius}(\kappa_0\circ g_z))$,
  since $\kappa_0\circ g_z$ equals the $z$-configuration step function a.e.
  By \cref{lem:error_vector_lipschitz} applied to the
  weights $\rho = \mathrm{radius}(\kappa\circ h_1\circ g_z)$ and $\rho_0 =
  \mathrm{radius}(\kappa_0\circ g_z)$,
  \[
    \lVert E_\kappa(z) - \mathrm{errorMap}(z)\rVert
      \le \int_0^{2\pi}\Bigl\lvert \tfrac{1}{\kappa(h_1(g_z\theta))}
        - \tfrac{1}{\kappa_0(g_z\theta)}\Bigr\rvert\,d\theta .
  \]
  Both weights lie in $[0,M]$, and $\lvert 1/u-1/v\rvert\le M^2\lvert u-v\rvert$
  for $u,v\ge 1/M$, so the integrand is $\le M^2\lvert\kappa\circ h_1 - \kappa_0\rvert\circ g_z$.
  Substituting $\phi = g_z(\theta)$ (a piecewise-affine change of variables with
  $g_z'\ge 1/K$) pulls the bound back to the $\phi$-line: the set where
  $\lvert\kappa(h_1\phi)-\kappa_0(\phi)\rvert>\varepsilon$ has measure
  $<\varepsilon$ (\cref{lem:exists_preliminary_reparam}), and its $g_z$-preimage
  has measure $<K\varepsilon$; off it the integrand is $\le M^2\varepsilon$, and on
  it the integrand is $\le 2M\cdot M = 2M^2$. Hence the integral is
  $\le M^2\varepsilon\cdot 2\pi + 2M^2\cdot K\varepsilon \le K' M^2\varepsilon$ for
  an absolute constant $K'$. Absorb $2\pi$ and the factor into $K$.
\end{proof}

\begin{lemma}

[Existence of a preconditioning reparametrization]
  \label{lem:exists_reparam_close}
  % Lean (private, not pinned): exists_reparam_kappaErrorMap_close, curvature_bounds, recip_diff_abs_le, integrableOn_of_measurable_bounded, measurable_stepCurvature_canonical
  \uses{def:kappa_error_map, def:error_map, lem:kappa_error_map_close,
        lem:exists_preliminary_reparam, lem:align_reparam_matches,
        lem:has_deriv_align_reparam, lem:align_reparam_change_of_var}
  Project-bespoke (the existence form of \cref{lem:kappa_error_map_close} that
  bundles the $\varepsilon$-choice). For any target margin $\mu>0$ there is an
  orientation-preserving circle reparametrization $h_1$ (\texttt{StrictMono},
  \texttt{Continuous}, quasi-periodic) such that
  $\lVert E_\kappa(z)-\mathrm{errorMap}(a,b,\delta)(z)\rVert<\mu$ for every
  $z\in\overline{\mathbb D}$. The same $h_1$ is the one of
  \cref{lem:exists_preliminary_reparam}, so it additionally carries the $C^1$
  regularity there: there is a continuous $v_1>0$ with $h_1'(\theta)=v_1(\theta)$
  for all $\theta$ (forwarded verbatim).
\end{lemma}
\begin{proof}
  \uses{lem:kappa_error_map_close, lem:exists_preliminary_reparam, lem:align_reparam_matches}
  By \cref{lem:kappa_error_map_close} the error-vector gap is bounded by the
  $L^1$ weight difference, so it suffices to make that integral $<\mu$ uniformly
  in $z$. Choose the tolerance $\varepsilon$ of \cref{lem:exists_preliminary_reparam}
  with $K M^2\varepsilon<\mu$ ($M=1/\min\kappa$, $K$ the slope bound of $g_z$),
  apply \cref{lem:exists_preliminary_reparam} to obtain $h_1$, and conclude by the
  change-of-variables / bad-set pullback through the slope-bounded $g_z$ (the
  $<\varepsilon$-measure bad set has $g_z$-preimage of measure $<K\varepsilon$; off
  it the integrand is $\le M^2\varepsilon$).
\end{proof}

\section{The reduction is justified}

% SOURCE: [DeTurck-Gluck], §6, p. 13, lines 243-248 (read from references/deturck-gluck-fourvertex.txt)
% SOURCE QUOTE: "But the point is that the above argument is robust, and hence
% applies equally well to the curvature function κ ◦ h1 which is ε-close in
% measure to κ0 , since the corresponding curves are C1-close. Hence there must
% also be a diffeomorphism h of the circle so that the curve with curvature κ ◦
% h1 ◦ h has error vector zero, and therefore closes up to form a strictly convex
% simple closed curve. Reparametrizing this curve as in section 4 shows that it
% realizes the preassigned curvature function κ , completing the proof of the
% converse to the Four Vertex Theorem for strictly positive curvature."
\begin{lemma}


\leanok
[Reduction is justified]
  \label{lem:reduction_justified}
  \lean{Gluck.reduction_justified}
  \uses{lem:abab_values, lem:exists_preliminary_reparam, def:kappa_error_map,
        lem:kappa_error_map_continuous, lem:kappa_error_map_close,
        lem:exists_reparam_close, lem:error_map_winds_boundary,
        thm:existence_of_zero}
  \textit{Source: DeTurck--Gluck, §6 (robustness of the winding argument).}
  Let $\kappa$ be a curvature function satisfying the four-vertex condition. Then
  there is a reparametrization $h$ of the circle (\texttt{StrictMono},
  \texttt{Continuous}, $h(\theta+2\pi)=h(\theta)+2\pi$) such that $\kappa\circ h$ is
  again a curvature function whose reconstruction weight $\rho = 1/(\kappa\circ h)$
  has $E(\rho)=0$; so by \cref{lem:reconstruct_periodic} the reconstruction curve
  $\alpha_\rho$ closes up. Moreover $h$ is a $C^1$ circle diffeomorphism with
  strictly positive derivative: there is a continuous $v>0$ with
  $h'(\theta)=v(\theta)$ for all $\theta$. This $C^1$ regularity is what lets the
  closed curve be reparametrized back by $h^{-1}$ to realize the original $\kappa$
  in \cref{thm:gluck_converse}.
\end{lemma}
\begin{proof}
  \uses{lem:abab_values, lem:exists_preliminary_reparam, lem:exists_reparam_close,
        lem:kappa_error_map_continuous, lem:kappa_error_map_close,
        lem:error_map_winds_boundary, thm:existence_of_zero,
        lem:has_deriv_align_reparam, lem:align_density_props}
        \leanok
  \emph{Step 1 (fix the winding margin).} Fix the configuration-disk radius
  $0<\delta\le\pi/8$ so that \cref{lem:error_map_winds_boundary} holds: the
  bicircle error map $E_{\mathrm{bi}} = \mathrm{errorMap}(a,b,\delta)$ is
  continuous, nonvanishing on $\partial\mathbb{D}$ with winding $\pm1$, and (from
  the explicit linear model, $\lVert V\rVert\approx\delta$) there is a margin
  $\mu := \tfrac12\min_{\partial\mathbb{D}}\lVert E_{\mathrm{bi}}\rVert > 0$.

  \emph{Step 2 (precondition $\kappa$).} Let $M=1/\min\kappa$ (finite, $\kappa>0$
  continuous on the compact circle). With $K$ the constant of
  \cref{lem:kappa_error_map_close}, choose $\varepsilon>0$ with $K M^2\varepsilon<\mu$
  and apply \cref{lem:exists_preliminary_reparam} (with the crossing data
  $0<a<b$, $c_1<\dots<c_4$ from \cref{lem:abab_values}) to obtain $h_1$.

  \emph{Step 3 (transfer the winding).} Form $E_\kappa$ (\cref{def:kappa_error_map}),
  continuous on $\overline{\mathbb{D}}$ (\cref{lem:kappa_error_map_continuous}). By
  \cref{lem:kappa_error_map_close}, $\lVert E_\kappa(z)-E_{\mathrm{bi}}(z)\rVert\le
  K M^2\varepsilon<\mu$ for all $z$, in particular on $\partial\mathbb{D}$. Hence
  on $\partial\mathbb{D}$, $\lVert E_\kappa\rVert\ge\lVert E_{\mathrm{bi}}\rVert-\mu
  \ge\mu>0$, so $E_\kappa\ne0$ there and the straight-line homotopy
  $E_{\mathrm{bi}}+\tau(E_\kappa-E_{\mathrm{bi}})$ is nonvanishing on
  $\partial\mathbb{D}$ for $\tau\in[0,1]$; winding is a homotopy invariant, so
  $E_\kappa$ also winds $\pm1\ne0$ on $\partial\mathbb{D}$.

  \emph{Step 4 (extract the zero).} By the planar degree theorem
  \cref{thm:existence_of_zero}, $E_\kappa$ has an interior zero $z_0$:
  $\mathrm{errorVector}(\mathrm{radius}(\kappa\circ h_1\circ g_{z_0}))=0$. Set
  $h = h_1\circ g_{z_0}$; then $\kappa\circ h$ is a curvature function (composition
  of a curvature function with orientation-preserving circle homeomorphisms) with
  $E(\mathrm{radius}(\kappa\circ h))=0$, which is the claim. (By
  \cref{lem:reconstruct_periodic} the reconstruction $\alpha_{1/(\kappa\circ h)}$ then closes
  up; reparametrizing back as in \cref{thm:gluck_converse} realizes $\kappa$.)

  \emph{$C^1$ regularity of $h$.} Both factors are $C^1$ with strictly positive
  derivative: $h_1'=v_1>0$ by \cref{lem:exists_reparam_close} (forwarding
  \cref{lem:exists_preliminary_reparam}), and $g_{z_0}'=w_{z_0}\ge 2/3>0$ by
  \cref{lem:has_deriv_align_reparam} with the density bounds of
  \cref{lem:align_density_props}. By the chain rule
  $h'(\theta)=v_1(g_{z_0}(\theta))\,w_{z_0}(\theta)=:v(\theta)$, a product of two
  continuous strictly positive functions, so $v$ is continuous and $v>0$ and
  $h\in C^1$ as claimed.

  \emph{The robustness is exactly $L^1$/measure continuity of $E$
  (\cref{lem:error_vector_lipschitz}), not $C^1$-closeness of curves.}
\end{proof}
